\documentclass[11pt,reqno]{amsart}
\usepackage[T1]{fontenc}
\usepackage{lmodern,amsmath,amssymb,mathtools,booktabs,array,longtable}
\usepackage[margin=1.05in]{geometry}
\usepackage{microtype}
\usepackage{xcolor,xr-hyper}
\usepackage[unicode,colorlinks=true,linkcolor=blue!45!black,citecolor=blue!45!black,urlcolor=blue!45!black]{hyperref}
\usepackage{enumitem,needspace,etoolbox}
\setcounter{tocdepth}{1}
\raggedbottom
\hypersetup{pdftitle={A determinant bound of 331 for the tenth discriminant},pdfauthor={Stephen Chung}}
\numberwithin{equation}{section}
\newtheorem{theorem}{Theorem}[section]
\newtheorem{proposition}[theorem]{Proposition}
\newtheorem{lemma}[theorem]{Lemma}
\newtheorem{corollary}[theorem]{Corollary}
\theoremstyle{remark}
\newtheorem{remark}[theorem]{Remark}
\newcommand{\OO}{\mathcal O}
\newcommand{\ZZ}{\mathbb Z}
\newcommand{\QQ}{\mathbb Q}
\newcommand{\CC}{\mathbb C}
\newcommand{\FF}{\mathbb F}
\newcommand{\NN}{\mathbb Z_{\geq0}}
\newcommand{\ord}{\operatorname{ord}}
\newcommand{\rank}{\operatorname{rank}}
\newcommand{\wt}{\operatorname{wt}}
\newcommand{\N}{\operatorname{N}}
\newcommand{\floor}[1]{\left\lfloor #1\right\rfloor}
\newcommand{\ceil}[1]{\left\lceil #1\right\rceil}
\newcommand{\pos}[1]{\left(#1\right)_+}
\newcommand{\fall}[2]{(#1)_{\underline{#2}}}
\newcommand{\rise}[2]{(#1)_{#2}}
\newcommand{\dd}{\,\mathrm d}
\newcommand{\code}[1]{\texttt{\detokenize{#1}}}
\newcommand{\arith}{\mathcal A}
\newcommand{\metric}{\mathcal U}
\newcommand{\supp}{\operatorname{supp}}
\setlength{\emergencystretch}{2em}
\title[A determinant bound for the tenth discriminant]{A determinant bound of 331 for the tenth discriminant}
\author[Stephen Chung]{Stephen Chung}
\date{Working draft, 28 October 2026}
\thanks{This work is based on an externally verified finding made by autonomous
rStation agents (\url{https://rstation.ai/}). See the section on
\hyperref[sec:ai-use]{AI use} for details.}
\keywords{Imaginary quadratic fields, class number one, interpolation determinants, Frobenius, explicit bounds}
\BeforeBeginEnvironment{lemma}{\Needspace{7\baselineskip}}
\BeforeBeginEnvironment{proposition}{\Needspace{9\baselineskip}}
\BeforeBeginEnvironment{theorem}{\Needspace{7\baselineskip}}
\makeatletter
\def\@setauthors{%
  \vskip 14pt
  \begingroup
  \centering\normalsize
  Stephen Chung\par
  \vskip 4pt
  {\small DualverseAI\par}%
  \endgroup
  \vskip 20pt
}
% Table captions are below the table; place their separation above them.
\setlength{\belowcaptionskip}{0pt}
\pretocmd{\@makecaption}{\ifnum\@tempcnta>63\relax\vskip8pt\fi}{}{}
\makeatother
\begin{document}
\begin{abstract}
We address the challenge of obtaining a small explicit upper bound for a
hypothetical tenth imaginary quadratic discriminant of class number one,
without using Heegner's methods or the known classification. We also make no quasi-Riemann-hypothesis assumption. We prove $D\leq331$ for every negative fundamental discriminant $-D$
with class number one, improving the earlier bound $D\leq498$. The earlier bound and elementary norm
arguments leave six possible discriminants. To exclude them, we sharpen the
rank criterion for an interpolation matrix by retaining the exact degree and
vanishing-order bounds in its proof. This permits the determinant comparison
to be applied with fewer evaluation points for each chosen polynomial space.
A count of norm representations supplies the auxiliary inert primes.
The arithmetic and spectral estimates are those used for the earlier bound;
the six new numerical comparisons use directed rounding.
\end{abstract}
\maketitle
\section{Introduction}\label{sec:intro}
Let $-D<0$ be a fundamental discriminant, and write $\OO_D$ for the
ring of integers of $\mathbb Q(\sqrt{-D})$.
Gauss's class number one problem asks when $\OO_D$ has unique factorization,
or equivalently when its ideal class number $h(-D)$ is one.
Heegner~\cite{heegner}, Baker~\cite{baker} and Stark~\cite{stark} proved that
this occurs exactly for
\begingroup
\abovedisplayskip=14pt
\belowdisplayskip=14pt
\abovedisplayshortskip=14pt
\belowdisplayshortskip=14pt
\[
 D\in\{3,4,7,8,11,19,43,67,163\}.
\]
\endgroup
Here we obtain a weaker bound by the determinant method, without using
Heegner's methods or the known classification.

\begin{theorem}\label{thm:main}
Let $-D<0$ be a fundamental discriminant with $h(-D)=1$.  Then
$D\leq331$.
\end{theorem}

The question is how small an explicit bound for a hypothetical tenth
class-number-one discriminant can be obtained from such an argument.  This
is the setting of the challenge~\cite{tao2026}, which uses the class
number one problem to test recent determinant arguments concerning
Landau--Siegel zeros~\cite{green2026,taosketch}.  Our earlier paper established
$D\leq498$ by comparing an upper bound for the norm of an interpolation
determinant with the prime powers forced to divide that norm~\cite{chung498}.
The present paper keeps that determinant and its spectral size estimate.
We sharpen the rank criterion by retaining the exact degree and vanishing
estimates in the earlier proof.  For a fixed polynomial space and fixed
derivative orders, this permits fewer evaluation points and improves the
upper bound for the determinant.

\subsection{Overview of the proof}\label{subsec:overview}
Suppose that $h(-D)=1$ and $D>331$.  Every element of $\OO_D$ has the form
$(a+b\sqrt{-D})/2$ with $a,b\in\mathbb Z$, so every non-rational element
has norm at least $D/4$.  If a prime $p$ splits or ramifies, an ideal of
norm $p$ must be principal.  Its generator cannot be a rational integer,
whose norm would be a square.  Thus every prime $p<D/4$ is inert.  In
particular, $D$ is odd, since otherwise $2$ ramifies.  Nor can $D$ be
composite: its least prime divisor would ramify and be at most
$\sqrt D<D/4$.  Since $-D$ is fundamental, the resulting odd prime $D$ is
congruent to $3$ modulo $4$; it cannot be $7$ modulo $8$, since then $2$
splits.  Hence $D$ is prime with $D\equiv3\pmod8$.

The bound $D\leq498$ from~\cite{chung498} now leaves only
\[
  D\in\{347,379,419,443,467,491\}.
\]
We exclude these six values using the determinant estimates
of~\cite{chung498} and the sharper rank criterion proved below.

We form the matrix by evaluating monomials $z^rw^s$ with $r+Ks<R$,
and their normalized $z$-derivatives, at pairs $(z,\bar z)$ with
$z\in\OO_D$.  The parameter $K$ determines the shape of this triangular
set of exponents.  At an inert prime $p$, conjugation agrees with
Frobenius, so $z^p-\bar z$ is divisible by $p$.  Expressing the monomials
in terms of powers of $z^p-w$ therefore gives prime powers dividing every
maximal minor, with a loss depending on the selected derivative orders.
If one of these minors is nonzero, its norm is a positive
integer with this prescribed divisibility.  On the other hand, choosing
the evaluation points near the origin gives an Archimedean upper bound
for the same norm.  We obtain a contradiction when the forced divisor
exceeds that upper bound.  The arithmetic estimate also uses larger inert primes: class number one
lets us bound the loss from non-inert primes by counting norm values.
The upper bound keeps the derivative columns at each node together,
using the spectral estimates of~\cite{chung498} to account for their
correlations.

The rank criterion determines how few points can be used while ensuring
that a nonzero maximal minor exists.  If the polynomial space has dimension
$M$, and we take $J$ derivative orders at each of $n$ points, the necessary
dimension inequality is $nJ\geq M$.  The sufficient inequality stated
in~\cite{chung498} does not hold uniformly over the auxiliary-prime
intervals used here. Its proof, however, contains more precise information.

As in that proof, we rescale the polynomial rows and derivative columns
by known powers of an auxiliary
inert prime before reducing modulo the prime.  A relation among the reduced
rows would give a polynomial in two variables with prescribed local
vanishing; its second variable records the powers of $z^p-w$.  Applying a differential operator and taking a determinant of
coefficients then produces a nonzero polynomial in one variable.  We bound
its degree and the multiplicity of its zeros, retaining their dependence
on the degree of the original relation in its second variable.  These
more precise bounds show that the total multiplicity exceeds the degree,
even when the earlier sufficient inequality does not apply.  This proves
full row rank for the smaller node counts.

We still need an auxiliary inert prime for which this rank criterion
applies.  The required primes are larger than $D/4$, where the elementary
argument no longer guarantees inertness.  Our chosen intervals lie above
$D$, so their primes cannot ramify.  If a prime in one of these intervals
were split, class number one would represent it by the norm form
\[
 p=x^2+xy+\frac{D+1}{4}y^2.
\]
Counting the possible representations bounds the number of split primes
in the interval.  When $y$ is even, the
norm is congruent to $x^2$ modulo $2$; an odd prime norm therefore requires
$x$ odd, so the admissible values of $x$ are spaced by two.  We use this parity refinement
for the two intervals needed when $K=17$.  Since the count is smaller than
the number of primes in the interval, at least one prime there is inert.
We also place each interval above twice the radius of the disk containing
the nodes.  Distinct nodes then remain distinct after reduction at every
prime in the interval, as the rank criterion requires.

For $D=491,467,443,419$ we use $K=43$.  For $D=379$ and $D=347$, the
stronger rank criterion permits $K=17$. At fixed $R$, reducing $K$ allows more powers of the conjugate
variable, changing both the divisibility and the Gram-matrix estimate. Their combined
comparison is more favorable at these two discriminants.  For each chosen polynomial space, the sharper criterion brings the
number of points close to the minimum $M/J$. With these
choices, the arithmetic lower bound exceeds the Meixner--Schur upper bound
for each of the six candidates.  All six values are therefore excluded,
proving Theorem~\ref{thm:main}.

The rest of the paper is organized around these steps.  Section~\ref{sec:comparison}
recalls the determinant inequality and states the imported size estimate.
Section~\ref{sec:rank} proves the exact rank criterion.
Section~\ref{sec:primes} supplies the auxiliary inert primes and fixes
the three parameter families, while Section~\ref{sec:completion} gives
the six numerical comparisons.  The accompanying calculations
and source code are available at
\url{https://github.com/dualverse-ai/tenth_discriminant}.

\section{The determinant comparison}\label{sec:comparison}
We recall the determinant and the estimates from~\cite{chung498} that will
be used. The comparison applies to every nonzero maximal minor. Thus the
rank argument in the next section need only establish that such a minor
exists; it need not identify its columns.

\subsection{The matrix and its normalization}
Assume that $h(-D)=1$ and $D>331$. Write
$\omega=(1+\sqrt{-D})/2$ and $\OO_D=\mathbb Z[\omega]$.
Choose integers $K\geq2$, $S\geq1$ and $1\leq J\leq R:=KS$, and put
\begin{equation}\label{eq:triangle}
 \Lambda=\{(r,s)\in\mathbb Z_{\geq0}^2:r+Ks<R\},\qquad
 M=|\Lambda|=\frac{R(S+1)}2.
\end{equation}
Let $z_1,\ldots,z_n$ be the $n$ elements of $\OO_D$ nearest to the origin,
with ties broken arbitrarily. We use the Hasse derivatives in $z$, with the
second variable $w$ held fixed: $\partial_z^{[j]}P$ is the coefficient of
$t^j$ in $P(z+t,w)$. The matrix is
\begin{equation}\label{eq:jet-matrix}
 V_{(r,s),(i,j)}=\binom rjz_i^{r-j}\bar z_i^{\,s},
 \qquad (r,s)\in\Lambda,\quad 1\leq i\leq n,\quad 0\leq j<J.
\end{equation}
The entry is zero if $j>r$. Dividing ordinary derivatives by $j!$ removes their factorial
factors while retaining algebraic-integer entries. For any selection
$C$ of $M$ columns, let $\Delta_C=\det V_C$.

The nearest-node estimate of~\cite[Lemma 2.5]{chung498} gives
\begin{equation}\label{eq:node-disk}
 |z_i|\leq\rho_i:=\sqrt{\frac{i\sqrt D}{2\pi}}+\delta,
 \qquad \delta=\frac{\sqrt{D+4}}4.
\end{equation}
The lattice has covolume $\sqrt D/2$ and covering radius
at most $\delta$: its horizontal rows have vertical spacing $\sqrt D/2$,
and the points on each row have spacing one. The disk of radius
$\rho-\delta$ is therefore covered by Voronoi cells of lattice points in
the disk of radius $\rho$. Comparing areas shows that this latter disk
contains at least $2\pi(\rho-\delta)^2/\sqrt D$ lattice points, proving
\eqref{eq:node-disk}.

Put $\rho=\rho_n$, $\theta=M/(nJ)$, and
\begin{equation}\label{eq:normalization}
 W=\frac{R^2S}{3},\qquad W_1=2\sum_\Lambda r,\qquad W_2=2\sum_\Lambda s.
\end{equation}
The ratio $\theta$ is the proportion of available columns used by a maximal
minor. In particular, $\theta\leq1$ is necessary for full row rank. Direct
summation over the triangle gives
\begin{equation}\label{eq:radius-coefficient}
 \frac{W_1+W_2}{W}
 =1+\frac1K+\frac{3(K-1)}{2KS}+\frac{K-5}{2KS^2}.
\end{equation}

\subsection{The arithmetic estimate}\label{subsec:primes-import}
We use the local divisibility formulas of the earlier paper, with all
quantities evaluated at the present parameters. Their general derivation is
given in Appendices~A and~B of~\cite{chung498}.
We record the hypotheses and the meaning of the resulting terms, since the
finite comparisons in that paper use different parameters.

Suppose throughout this subsection that $D>163$ is prime,
$D\equiv3\pmod8$, and $h(-D)=1$. Put $m_D=(D+1)/4$.
Let $P$ be the prime cutoff in the arithmetic sums. The parameters used
here satisfy
\begin{equation}\label{eq:arithmetic-hypotheses}
 \begin{gathered}
 2\leq K<D/4,\quad R=KS,\quad 1\leq J\leq R,\quad M\leq nJ,\\
 P\geq\max(S,J-1),\qquad
 F_{\max}:=\lfloor P/m_D\rfloor<m_D^4.
 \end{gathered}
\end{equation}
The local divisibility estimates give a number $\arith$ and a derivative
loss $b(j)$ such that every nonzero maximal minor satisfies
\begin{equation}\label{eq:arithmetic-interface}
 W\arith-2\sum_C b(j)\leq\log|\det V_C|^2.
\end{equation}
We recall their construction to specify how the estimates apply at the
present parameters. An element is \emph{primitive} if it is not divisible
in $\OO_D$ by a rational integer greater than one.

All primes below $K$ are then inert. The small-prime estimates apply without
exception, and the cutoff $P$ includes the derivative debits, the ramified
prime, and the primes contributing to the conjugate-variable estimates.
Summing the local row bounds gives
\begin{equation}\label{eq:arithmetic-parts}
 \arith=\arith_0-L_{\rm norm}+C_2+C_3.
\end{equation}
Here $\arith_0$ provisionally treats primes up to $P$ as inert.
Class number one represents every non-inert prime as the norm of an
algebraic integer. Summing over primitive norm representatives therefore
gives an upper bound $L_{\rm norm}$ for the necessary correction, retaining
the split and ramified local contributions. This sum also includes composite
norm values. The nonnegative terms $C_2,C_3$ correct this overestimate for
products of two or three prime factors, with the representation
multiplicities supplied by unique factorization. Omitting other composite
corrections only weakens the lower bound.

Some split-prime column contributions are retained individually. Let
$\mathcal P_s(D)$ be the split primes at most $q_{\rm cap}$, where the
cutoff satisfies
\begin{equation}\label{eq:slope-cutoff}
 q_{\rm cap}\leq\min(J-1,F_{\max}m_D,m_D^4-1).
\end{equation}
Define the remaining derivative loss and its factorial remainder by
\begin{equation}\label{eq:arithmetic-debit}
 \begin{split}
 b(j)&=\sum_{q\leq j}\lfloor j/q\rfloor\log q
       -\sum_{q\in\mathcal P_s(D)}\lfloor j/q\rfloor\log q,\\
 F(j)&=\log(j!)-b(j)
       =\sum_q\sum_{a\geq2}\lfloor j/q^a\rfloor\log q
        +\sum_{q\in\mathcal P_s(D)}\lfloor j/q\rfloor\log q.
 \end{split}
\end{equation}
Sums indexed by $q$ are over primes. The split-prime contributions outside this
window are bounded uniformly over column selections and already included
in \eqref{eq:arithmetic-parts}.

The prime sums in the window are bounded using norm representatives and
composite products. Every prime factor of a primitive norm is at least
$m_D$. Thus a composite primitive norm in \eqref{eq:slope-cutoff} has at most
three prime factors, counted with multiplicity, and each factor is at
most $F_{\max}$. To obtain the factor list up to $F_{\max}$, the finite enumeration
first removes reducible primitive elements. A reducible
primitive element of norm at most $F_{\max}<m_D^4$ has an irreducible factor
of norm at most $\sqrt{F_{\max}}<m_D^2$, and a remaining factor of norm at
most $F_{\max}/m_D$. Appendix~B of~\cite{chung498} proves that these ranges construct the complete factor list. Consequently the subtraction includes every
composite in the window, even for $P>m_D^4$.

Using Appendix~B of~\cite{chung498}, we choose a nonnegative lower bound $\sigma$ for
\[
 \sum_{\substack{q,\ a\geq2\\q^a\leq J-1}}\frac{\log q}{q^a}
       +\sum_{q\in\mathcal P_s(D)}\frac{\log q}{q},
\]
and upper bounds $h_j$ for $H(j)=j\sigma-F(j)$.
These bounds determine $c$ by the capacity calculation below.
The explicit finite weights, slope bounds and remainder estimates are
those in Appendix~B of~\cite{chung498}. The hypotheses above hold for all three parameter
families below.


\subsection{Combining the derivative terms}
Legendre's formula writes $\log(j!)$ as
$\sum_q\sum_{a\geq1}\lfloor j/q^a\rfloor\log q$. Its first-power
terms cancel the corresponding derivative loss in $b(j)$, leaving the
higher prime powers and split-prime contributions in $F(j)$. We use the
remainder $H(j)=j\sigma-F(j)$ defined above,
and choose a constant $c$ satisfying
\begin{equation}\label{eq:column-remainder}
 2\sum_{(i,j)\in C}H(j)\leq Mc\qquad(|C|=M).
\end{equation}
This bound uses the fact that each derivative order is available at only
$n$ nodes. Specifically, if $h_{(1)}\geq\cdots\geq h_{(J)}$ are upper
bounds for $H(0),\ldots,H(J-1)$ arranged in decreasing order, and
$M=nt+r$ with $0\leq r<n$, one may take
\[
 Mc=2n\sum_{a=1}^{t}h_{(a)}+2r h_{(t+1)}.
\]
The last term is omitted when $r=0$. The formulas in Appendix~B of~\cite{chung498} construct $\sigma$ and these upper bounds from prime powers and norm
values, so the calculation does not require identifying the split primes
at the candidate discriminant.

Set $u_i=z_i/\rho$ and $\ell=\log(R/\rho)-\sigma$. Define the normalized
matrix
\begin{equation}\label{eq:residual-matrix}
 B_{(r,s),(i,j)}(\ell)
 =\frac{\fall rj}{R^j}e^{\ell j}u_i^{r-j}\bar u_i^{\,s},
\end{equation}
where $\fall rj=r(r-1)\cdots(r-j+1)$. Exact diagonal scaling gives
\[
 B_C(\ell)=\operatorname{diag}_{\Lambda}(\rho^{-r-s})\,
 V_C\,\operatorname{diag}_C(j!e^{-\sigma j}).
\]
Taking squared absolute determinants and using \eqref{eq:arithmetic-interface} and
\eqref{eq:column-remainder}, we obtain
\begin{equation}\label{eq:comparison-before-size}
 W\arith\leq (W_1+W_2)\log\rho+
                \log|\det B_C(\ell)|^2+Mc.
\end{equation}
Since $e^{\ell j}$ scales column $(i,j)$ and $j\geq0$, replacing
$\ell$ by an upper bound only increases $|\det B_C(\ell)|^2$. It
therefore suffices to bound the normalized determinant uniformly over $C$.

\subsection{The imported size estimate}
We use the Meixner and Schur estimates from~\cite[Section~4]{chung498}.
Their role is to retain the correlations between derivative columns at
one node. We recall the quantities entering the bound so that its
application at the new parameters is explicit.

Choose $a>0$, $0<v<1$, $\beta<0$, and put $b=aR$. Weight row $(r,s)$ by
$e^{-\lambda_{r,s}/2}$, where
\[
 e^{-\lambda_{r,s}}=\frac{(b)_r}{r!}v^r e^{-\beta s},
 \qquad (b)_r=b(b+1)\cdots(b+r-1).
\]
Let $Y=\operatorname{diag}(e^{-\lambda_{r,s}/2})B(\ell)$, and let $G_i=Y_i^*Y_i$
be the Gram matrix of its $J$ columns at node $i$. Fischer's inequality gives
\begin{equation}\label{eq:fischer}
 \log|\det B_C(\ell)|^2
 \leq\sum_\Lambda\lambda_{r,s}
       +\sum_i\log\det G_i[C_i,C_i],
\end{equation}
where $C_i$ is the set of orders selected at that node. By~\cite[Lemma~4.1]{chung498}, each principal minor
of $G_i$ is a polynomial with nonnegative coefficients in $q=|u_i|^2$;
therefore it increases when the node is moved radially outward.

Write $x_+=\max(0,x)$. At an upper squared radius $q>0$, set
\[
 d_q=(\log q-\beta)_+,\qquad v_q=ve^{-d_q/K},\qquad
 w=v_qq,\qquad \mathcal B_q=S e^{Sd_q}.
\]
After a diagonal unitary conjugation removing the node angle, the Gram
matrix is bounded in the positive semidefinite order by
\begin{equation}\label{eq:moment-matrix}
 \mathcal G_{jk}(q)=\mathcal B_q
 \left(\frac{e^\ell}{R\sqrt q}\right)^{j+k}
 \sum_{r\geq0}\frac{(b)_r}{r!}w^r\fall rj\fall rk,
 \qquad 0\leq j,k<J.
\end{equation}
The conjugation preserves every principal minor. Here $0<w<1$ for the
parameters below. The negative-binomial moments in
this formula have an explicit factorization in Meixner polynomials. This
supplies two bounds for an arbitrary principal minor: one uses the largest
singular values of an affine composition operator, and the other its
complementary singular product. A third bound sums all principal minors of
the prescribed size, using the Schur-function Cauchy identity. These are
respectively equations (4.10), (4.11), and (4.12) of~\cite{chung498}.
Each holds for every subset of derivative orders; no particular subset is
assumed to maximize the determinant.

Let $F_q(m)$ be the minimum of these three upper bounds for the logarithm
of an $m$-column principal minor, allowing the finite parameter choices in
Appendix~C of~\cite{chung498}. For a chosen $a_0>0$, take
$\Phi(q)\geq0$ such that
\begin{equation}\label{eq:H-occupancy}
 F_q(m)\leq a_0Rm+R^2\Phi(q)\qquad(0\leq m\leq J).
\end{equation}
We divide the range of $m/R$ into $N_m$ equal intervals and establish
the bound on each whole interval before minimizing over estimates.
For the singular-value bounds, the sum of the $m$ largest column prices
is bounded using $N_z$ equal intervals of $j/R\in[0,\varepsilon]$,
where $\varepsilon=J/R$. Appendix~C of~\cite{chung498} bounds the price and the number of orders in each
interval. Summing \eqref{eq:H-occupancy} in
\eqref{eq:fischer} uses $\sum_i|C_i|=M$, rather than allowing every node to
choose its worst number of columns independently.

To bound the radii, put
\[
 \eta=\frac{\delta}{\sqrt{n\sqrt D/(2\pi)}},\qquad
 Q(t)=\left(\frac{\sqrt t+\eta}{1+\eta}\right)^2,
 \qquad \varepsilon=J/R.
\]
Then $|u_i|^2\leq Q(i/n)$. Divide $(0,1]$ into $N_t$ equal intervals and
move the nodes in the $k$th interval to squared radius $Q(k/N_t)$.
Each interval contains at most $n/N_t+1$ nodes. The resulting bound is
\begin{equation}\label{eq:metric-mesh}
 U_{\rm var}=U_{\rm row}+\frac{3a_0(S+1)}{2S}
 +3\left(\frac{1+1/S}{2\theta\varepsilon N_t}+\frac1S\right)
       \sum_{k=1}^{N_t}\Phi(Q(k/N_t)),
\end{equation}
where $U_{\rm row}$ is an upper bound for
$W^{-1}\sum_\Lambda\lambda_{r,s}$. In particular,
$W^{-1}\log|\det B_C(\ell)|^2\leq U_{\rm var}$ for every $C$.
For the finite formulas in Appendix~C of~\cite{chung498},
we require $R>3$, $aR>1$, $0<a\leq1$, $0<\varepsilon\leq1$,
$\varepsilon(a+2\varepsilon)\leq1$, and
\[
 -(a+1)\log(a+1)+a\log a-\log v<0.
\]
The last condition gives the signed row-sum bound. For the Schur bound at
radius $q$, put
\[
 \zeta=\ell+\tfrac12\log v_q-\log(1-w),\qquad z_*=a^{-2}e^{-2\zeta}.
\]
Its finite evaluation is used when $z_*<1$ and
$\varepsilon(a+\varepsilon)z_*<a^2(1-z_*)^2$. These are respectively the
moment and concavity conditions in Appendix~C of~\cite{chung498}. The other
two bounds remain available wherever the latter conditions do not hold.

\begin{proposition}[Determinant comparison]\label{prop:comparison}
Assume the arithmetic estimates \eqref{eq:arithmetic-interface} and
\eqref{eq:column-remainder}, and the size estimate \eqref{eq:metric-mesh}.
If $V$ has full row rank, then
\begin{equation}\label{eq:master}
 \arith\leq\metric:=\frac{W_1+W_2}{W}\log\rho+U_{\rm var}
                       +\frac{Mc+2M\log M}{W}.
\end{equation}
\end{proposition}
\begin{proof}
Full row rank supplies a nonzero minor. Apply
\eqref{eq:comparison-before-size} and the uniform size bound to that minor.
The nonnegative term $2M\log M/W$ retains the earlier numerical
normalization~\cite{chung498}.
\end{proof}
Thus a positive lower bound for $G:=\arith-\metric$ excludes the proposed
field. The arithmetic and size estimates depend on the actual node count
$n$ throughout. In particular, reducing $n$ changes both the disk radius
and the capacity bound for the derivative orders.

\section{An exact rank criterion}\label{sec:rank}
The determinant comparison requires the matrix $V$ to have full row rank.
The necessary dimension inequality $nJ\geq M$ does not by itself establish
this. We use the same reduction at an inert prime as in the earlier bound,
but retain the individual degree and multiplicity estimates throughout the
argument. This gives a stronger sufficient criterion at the parameters needed
here.

The obstruction to direct reduction is the identity
$\bar z=z^p$ modulo an inert prime $p$. Thus $z^p-w$ vanishes at every
reduced node, and some of the rows become dependent. Dividing out known
powers of $p$ before reducing preserves more information. A relation among
the resulting rows would produce a polynomial with prescribed vanishing
along several curves. We will construct from it a nonzero polynomial in
one variable, and show that its zeros have total multiplicity greater than
its degree.

\begin{theorem}[Exact rank criterion]\label{thm:exact-rank}
Let $K\geq2$ and $S\geq1$ be integers, put $R=KS$, and let $1\leq J\leq R$.
Suppose that $p$ is inert in $\OO_D$, that
\[
 K<p<R\leq p^2,
\]
and that the nodes $z_1,\ldots,z_n\in\OO_D$ have distinct residues modulo
$p\OO_D$.
Set $d=\lfloor(R-1)/p\rfloor$, and define
\begin{equation}\label{eq:rank-caps}
 \begin{split}
 D_m&=p\left\lfloor\frac{R-mp-1}{K}\right\rfloor
             +\bigl((R-mp-1)\bmod K\bigr)\qquad(0\leq m\leq d),\\
 B_e&=\sum_{m=0}^e(D_m-m)\qquad(0\leq e\leq d).
 \end{split}
\end{equation}
For $l\geq0$ put $L_l=\min(p,\max(0,J-lp))$, and let
\begin{equation}\label{eq:rank-multiplicity}
 \mu_e=\sum_{l=0}^e\max(0,L_l-e+l)\qquad(0\leq e\leq d).
\end{equation}
If
\begin{equation}\label{eq:exact-rank-ineq}
 n\mu_e>B_e\qquad(0\leq e\leq d),
\end{equation}
then the matrix \eqref{eq:jet-matrix} formed with these nodes has full
row rank.
\end{theorem}

Here $D_m$ is the largest degree of a coefficient polynomial arising from
layer $m$ of the reduced relation. The quantity $B_e$ bounds the degree of
the one-variable polynomial constructed from a relation of degree $e$ in
the second variable, while $\mu_e$ bounds its order of vanishing at each
node. Keeping these quantities separately for every $e$ avoids the losses
in the earlier sufficient criterion.

\subsection{A relation after reduction modulo an inert prime}
Put $g=z^p-w$. The polynomials
\begin{equation}\label{eq:g-basis}
 g^m z^a w^s,\qquad
 0\leq m\leq d,\quad 0\leq a<p,\quad a+Ks<R-mp,
\end{equation}
form an integral basis of the space spanned by the triangular support
$z^rw^s$, $r+Ks<R$. Indeed, writing $r=mp+a$ identifies their leading
monomials with the original basis. Every other term replaces a factor
$z^p$ by $-w$, decreasing the weighted degree by $p-K>0$.
The change of basis is therefore unitriangular.

At node $i$ write
\[
 g(z_i+t,\bar z_i)=pU_i(t)+t^p,\qquad
 U_i(t)=\frac{(z_i+t)^p-t^p-\bar z_i}{p}\in\OO_D[t].
\]
The constant coefficient is integral because $p$ is inert, and the remaining
coefficients are integral because $p$ divides $\binom pv$ for $0<v<p$.
Divide each row in layer $m$ by $p^m$, and multiply each column of derivative
order $j=lp+v$, $0\leq v<p$, by $p^l$. Denote the resulting matrix by
$\widetilde V$. Its entries are still integral: a term using $t^{pq}$ in
the expansion of $(pU_i+t^p)^m$ can contribute to order $j$ only if $q\leq l$,
and its coefficient after rescaling contains $p^{l-q}$.
Only $q=l$ can survive reduction modulo $p$.

Write $k=\OO_D/p\OO_D$, which is a field, and put
$\xi_i=z_i\bmod p$ and $u_i(z)=(U_i\bmod p)(z-\xi_i)$.
For $l\leq m$ the reduced entry is
\begin{equation}\label{eq:scaled-entry}
 \binom ml[t^v]\bigl(u_i(\xi_i+t)^{m-l}(\xi_i+t)^a\xi_i^{ps}\bigr);
\end{equation}
for $l>m$ it is zero. The basis change and the scalings are invertible in
characteristic zero. Consequently, full row rank of $\widetilde V\bmod p$
will imply full row rank of $V$. These scalings are used only for this
purpose; the size and divisibility estimates concern the original matrix.

Suppose that the reduced rows satisfy a nontrivial relation, with
coefficients $c_{m,a,s}\in k$. Define
\begin{equation}\label{eq:rank-F}
 F(z,u)=\sum_{m=0}^d H_m(z)u^m,\qquad
 H_m(z)=\sum_{a,s}c_{m,a,s}z^{a+ps}.
\end{equation}
The exponents $a+ps$ within a layer are distinct because $0\leq a<p$.
Thus $F\ne0$. Let $e=\deg_u F\leq d<p$, where the last inequality follows from
$R\leq p^2$.
For $t=R-mp-1$, the maximum of $a+ps$ subject to
$a+Ks\leq t$ and $0\leq a<p$ is attained at
$s=\lfloor t/K\rfloor$, $a=t\bmod K$, since $p>K$.
It follows that
\begin{equation}\label{eq:rank-degree}
 \deg H_m\leq D_m.
\end{equation}
We will also use $D_m-D_{m+1}\geq p$ for $m<d$.
To see this without manipulating the floors, take a monomial attaining
$D_{m+1}$ and increase $s$ by one. Its weighted degree increases by $K<p$,
so it is admissible in layer $m$, and its degree $a+ps$ increases by $p$.

There are $L_l$ derivative orders $j=lp+v$ with $0\leq j<J$.
Since $(\xi_i+t)^{ps}\equiv\xi_i^{ps}\pmod{t^p}$ in characteristic $p$,
we may replace $\xi_i^{ps}$ in \eqref{eq:scaled-entry} by
$(\xi_i+t)^{ps}$ when extracting coefficients of order less than $L_l$.
The row relation therefore gives
\begin{equation}\label{eq:local-hasse-vanish}
 \ord_{z=\xi_i}\bigl((\partial_u^{[l]}F)(z,u_i(z))\bigr)\geq L_l
 \qquad(1\leq i\leq n,\ 0\leq l\leq e).
\end{equation}
Here $\partial_u^{[l]}$ is the Hasse derivative, defined by
$\partial_u^{[l]}u^m=\binom ml u^{m-l}$.

\subsection{A nonzero polynomial of controlled degree}
Differentiating the identity defining $U_i$ before reducing modulo $p$ gives
\begin{equation}\label{eq:ui-prime}
 u_i'(z)=z^{p-1}-(z-\xi_i)^{p-1}.
\end{equation}
Thus differentiation along the different curves $u=u_i(z)$ agrees, up to
a term vanishing to order $p-1$ at the relevant node, with the common operator
\[
 \mathcal D=\partial_z+z^{p-1}\partial_u.
\]
On $k(z)[u]$ this operator satisfies
\begin{equation}\label{eq:Dp-identity}
 \mathcal D^p=-\partial_u.
\end{equation}
Indeed, the $p$th iterate of a derivation in characteristic $p$ is a
derivation. It sends $z$ to zero and $u$ to $(p-1)!=-1$, which determines
its action on $k(z)[u]$.

We claim that $F,\mathcal DF,\ldots,\mathcal D^eF$ are linearly independent
over $k(z)$. Otherwise, since $F\ne0$, let $1\leq r\leq e$ be the first index at
which a dependence occurs. The span of $F,\ldots,\mathcal D^{r-1}F$ is then stable
under $\mathcal D$, including when $\mathcal D$ differentiates rational
coefficients. By \eqref{eq:Dp-identity} it is also stable under $\partial_u$.
It must therefore contain $F,\partial_uF,\ldots,\partial_u^eF$.
These are independent: since $e<p$, they have distinct degrees
$e,e-1,\ldots,0$ in $u$, with nonzero leading coefficients.
This contradicts the dimension $r\leq e$.

Consequently the coefficient determinant
\begin{equation}\label{eq:WF-def}
 W_F(z)=\det\bigl([u^m]\mathcal D^jF\bigr)_{0\leq j,m\leq e}
\end{equation}
is a nonzero polynomial. The identity
\[
 [u^m]\mathcal DG=([u^m]G)' +(m+1)z^{p-1}[u^{m+1}]G
\]
and $D_m-D_{m+1}\geq p$ give, by induction on $j$,
\[
 \deg\bigl([u^m]\mathcal D^jF\bigr)\leq D_m-j.
\]
Terms with $m>e$ are zero, and a negative degree bound means the
corresponding entry is zero. Each determinant term uses every row and every
column exactly once, so
\begin{equation}\label{eq:WF-degree}
 \deg W_F\leq\sum_{m=0}^eD_m-\sum_{j=0}^ej=B_e.
\end{equation}

\subsection{The multiplicity at each node}
Fix $i$. Expressing each row polynomial of \eqref{eq:WF-def} in powers of
$u-u_i(z)$ makes a unitriangular change of columns. It preserves the
determinant and replaces its entries by
\[
 T^{(i)}_{j,l}(z)=(\partial_u^{[l]}\mathcal D^jF)(z,u_i(z)).
\]
We first show that
\begin{equation}\label{eq:local-entry-order}
 \ord_{\xi_i}T^{(i)}_{j,l}\geq\max(0,L_l-j).
\end{equation}
Let $\mathcal D_i=\partial_z+u_i'(z)\partial_u$.
By \eqref{eq:ui-prime},
$\mathcal D-\mathcal D_i=(z-\xi_i)^{p-1}\partial_u$.
Repeated application lowers the order of this error by at most one at each
step; hence, for every polynomial $G$ and $0\leq j<p$,
\[
 \mathcal D^jG-\mathcal D_i^jG
       \in(z-\xi_i)^{p-j}k[z,u].
\]
Both operators commute with the Hasse derivatives in $u$.
After substituting $u=u_i(z)$, the entry $T^{(i)}_{j,l}$ therefore agrees,
modulo $(z-\xi_i)^{p-j}$, with the $j$th ordinary derivative of
$(\partial_u^{[l]}F)(z,u_i(z))$.
Equation~\eqref{eq:local-hasse-vanish} bounds its order from below by $L_l-j$ when
this is positive. Since $L_l\leq p$, the error has at least the same
order, proving \eqref{eq:local-entry-order}.

To use all these bounds together, consider a term in the determinant
expansion. Its order is at least
$\sum_{l=0}^e(L_l-j_l)_+$, where the $j_l$ are a permutation of
$0,\ldots,e$. For $a\geq b$ and $x\geq y$ one has
\[
 (a-x)_++(b-y)_+\leq(a-y)_++(b-x)_+.
\]
Thus pairing the largest $L_l$ with the largest row index minimizes this
sum. Since $L_0\geq\cdots\geq L_e$, its minimum is
$\sum_{l=0}^e(L_l-e+l)_+=\mu_e$.
Every determinant term has at least this order, and so does their sum:
\[
 \ord_{\xi_i}W_F\geq\mu_e.
\]
Since $W_F$ is nonzero and the residues are distinct, \eqref{eq:WF-degree} gives
\begingroup\postdisplaypenalty=10000
\[
 n\mu_e\leq\deg W_F\leq B_e,
\]
\endgroup
contrary to \eqref{eq:exact-rank-ineq}. This proves Theorem~\ref{thm:exact-rank}.

\section{Auxiliary primes and the choice of parameters}\label{sec:primes}
The rank theorem requires an inert prime large enough to distinguish the
nodes, but smaller than $R$ and with $R\leq p^2$. We obtain it by counting
norm representations in a short interval. This uses the class number one
hypothesis in the same way as the arithmetic estimate: every split prime
has to be the norm of an algebraic integer.

\begin{lemma}[An inert prime in an interval]\label{lem:prime-window}
Suppose that $D>331$, $h(-D)=1$, and $D<A<B$. For $y\geq1$, put
$c_y=1$ when $y$ is odd and $c_y=1/2$ when $y$ is even. Define
\begin{equation}\label{eq:split-count}
 E_D(A,B)=\sum_{y=1}^{\lfloor\sqrt{4B/D}\rfloor}
 \left[1+c_y\left(\sqrt{B-Dy^2/4}
              -\sqrt{\max(0,A-Dy^2/4)}\right)\right].
\end{equation}
If $E_D(A,B)<\pi(B)-\pi(A)$, then $(A,B]$ contains an inert prime.
\end{lemma}
\begin{proof}
The elementary reduction makes $D$ prime, so every prime in $(A,B]$ is
unramified. If such a prime $p$ splits, a principal ideal of norm $p$
gives an element $x+y\omega$ with
\[
 p=x^2+xy+\frac{D+1}{4}y^2
   =\left(x+\frac y2\right)^2+\frac D4y^2.
\]
Here $y\ne0$, since a rational integer has square norm. Changing the sign
and taking the conjugate if necessary, we may assume $y>0$ and
$x+y/2\geq0$. Therefore $y\leq\sqrt{4B/D}$, and, for a fixed $y$,
$x+y/2$ lies between the two square roots in \eqref{eq:split-count}.
An interval of length $L$ contains at most $1+L$ members of any translate
of $\mathbb Z$.

When $y$ is even, we can improve this count. Writing $y=2k$ gives
\[
 x^2+xy+\frac{D+1}{4}y^2\equiv x^2\pmod2.
\]
Since $p$ is odd, $x$ is odd. Its permitted values are consequently spaced
by two, and an interval of length $L$ contains at most $1+L/2$ of them.
Summing these bounds gives at most $E_D(A,B)$ split primes. If this is
smaller than the number of primes in the interval, one of them must be inert.
\end{proof}

\subsection{The three parameter choices}
The six discriminants require three choices of polynomial space and node
count, given in Table~\ref{tab:parameters}. Throughout we set
\begin{equation}\label{eq:parameter-formulas}
 S=\left\lfloor\frac{6P}{5K}\right\rfloor,\qquad R=KS,\qquad
 J=\lfloor\varepsilon_0R\rfloor.
\end{equation}
Here $P$ is the prime cutoff in the arithmetic estimate, and
$\varepsilon_0$ is a rational number; the actual ratio in the size estimate
is $\varepsilon=J/R$.
\begin{table}[ht]
\centering
\begin{tabular}{lrrrr}
\toprule
$D$ & $P$ & $K$ & $\varepsilon_0$ & $n$\\
\midrule
$419,443,467,491$ & $312\,000\,000$ & $43$ & $1/8$ & $35\,155\,200$\\
$379$ & $312\,000\,000$ & $17$ & $16/125$ & $86\,265\,600$\\
$347$ & $756\,000\,000$ & $17$ & $16/125$ & $208\,919\,200$\\
\bottomrule
\end{tabular}
\caption{Parameters for the six determinant comparisons. Each row specifies
$S,R,J$ through \eqref{eq:parameter-formulas}.}
\label{tab:parameters}
\end{table}

The ratios $\theta=M/(nJ)$ in the three rows are approximately
$0.990690$, $0.997262$, and $0.997782$. Thus these choices use only about
$0.94\%$, $0.27\%$, and $0.22\%$ more nodes than the necessary minimum
$M/J$. The exact rank criterion is needed at these counts: the coarser
sufficient inequality in~\cite[Theorem 5.1]{chung498} does not hold
uniformly over the corresponding auxiliary-prime windows.

Table~\ref{tab:prime-windows} gives the intervals used to supply an inert
prime. For the first row we use the weaker count obtained by setting all
$c_y=1$; it already suffices. Its displayed split-prime bound is the largest
of the four separately evaluated counts, not a count on the whole
interval $419\leq D\leq491$.
\begin{table}[ht]
\centering
\begin{tabular}{lrrrr}
\toprule
$D$ & $(A,B]$ & $\pi(B)-\pi(A)$ & Split-prime bound & $2\rho_n<$\\
\midrule
$419,443,467,491$ & $(200\,000,228\,000]$ & $2276$ & $2181.456$ & $22\,281$\\
$379$ & $(500\,000,515\,000]$ & $1135$ & $990.602$ & $32\,708$\\
$347$ & $(800\,000,820\,000]$ & $1465$ & $1366.923$ & $49\,785$\\
\bottomrule
\end{tabular}
\caption{Upper bounds for the number of split primes and the node diameter.
All decimal bounds are rounded upward.}
\label{tab:prime-windows}
\end{table}

In each row $D<A$, and the split-prime bound is smaller than the prime
count. Lemma~\ref{lem:prime-window} therefore supplies an inert prime
$p\in(A,B]$. The diameter bound shows that the nodes remain distinct
modulo $p\OO_D$: a nonzero element of that ideal has modulus at least $p$,
whereas a difference of two nodes has modulus less than $A$.
Moreover $K<A<B<R<A^2$, so $K<p<R<p^2$.

The remaining conditions in Theorem~\ref{thm:exact-rank} are the finite
integer inequalities $n\mu_e(p)>B_e(p)$. They hold for every prime in the
corresponding window and every $0\leq e\leq\lfloor(R-1)/p\rfloor$.
The least differences $n\mu_e(p)-B_e(p)$ for the three rows are respectively
\[
 13\,425\,248\,140,\qquad 6\,204\,049\,940,\qquad 89\,748\,485\,986.
\]
The sums defining these integers are evaluated directly; their complete
ranges and a short integer-arithmetic program are in the repository.
Thus the particular inert prime supplied by the count satisfies every
hypothesis of the rank theorem, and $V$ has full row rank in all six cases.

\section{Completion of the proof}\label{sec:completion}
It remains to evaluate the determinant comparison at the parameters of
Table~\ref{tab:parameters}. For the size estimate we use
\begin{equation}\label{eq:spectral-parameters}
 (a,v,\beta,a_0)=\left(\frac{39}{40},\frac{19}{50},-1,\frac1{200}\right),
 \qquad (N_t,N_z,N_m)=(128,256,512).
\end{equation}
The last two integers are the mesh sizes for the column-order capacity
bound and the intervals of possible $m/R$. The finite choices of affine
composition parameters are the same as in Appendix~C of~\cite{chung498}. These values satisfy the finite hypotheses of
Section~\ref{sec:comparison}. At each upper radial endpoint, the
calculation checks the moment and concavity conditions before using the
Schur bound.


For the arithmetic estimate we retain only the two- and three-factor
composite corrections. The cutoff for the derivative slope is
\[
 q_{\rm cap}=\min\{J-1,\,m_D^4-1,\,m_D\lfloor P/m_D\rfloor\},
 \qquad m_D=(D+1)/4.
\]
It is $46\,799\,995$ in the first parameter family, $47\,923\,198$ at
$D=379$, and $57\,289\,760$ at $D=347$. In each case the factor bounds
and column capacities in the arithmetic estimate are satisfied.

Table~\ref{tab:gaps} gives rigorous enclosures for the difference
$G=\arith-\metric$ between the arithmetic lower bound and the size upper
bound in Proposition~\ref{prop:comparison}. The calculations use directed
rounding, and the displayed endpoints are rounded outward.
\begin{table}[ht]
\centering
\begin{tabular}{rr}
\toprule
$D$ & Enclosure for $G$\\
\midrule
$347$ & $[0.00146083,\;0.00147313]$\\
$379$ & $[0.01186524,\;0.01187031]$\\
$419$ & $[0.00417435,\;0.00417636]$\\
$443$ & $[0.04790491,\;0.04790692]$\\
$467$ & $[0.11463345,\;0.11463546]$\\
$491$ & $[0.17810049,\;0.17810250]$\\
\bottomrule
\end{tabular}
\caption{The six positive comparison margins, with two- and three-factor
composite corrections.}
\label{tab:gaps}
\end{table}

\begin{proof}[Proof of Theorem~\ref{thm:main}]
Suppose $h(-D)=1$ and $D>331$. The bound $D\leq498$ from~\cite{chung498}
and the elementary reduction in Section~\ref{sec:intro} restrict $D$ to
$347,379,419,443,467,491$. Section~\ref{sec:primes} supplies the
auxiliary inert prime and verifies the hypotheses of
Theorem~\ref{thm:exact-rank}, so the interpolation matrix has a nonzero
maximal minor in each case. Proposition~\ref{prop:comparison} requires
$G\leq0$, whereas every lower endpoint in Table~\ref{tab:gaps} is positive.
This excludes all six possibilities.
\end{proof}

The bound $331$ shows that a sharper nonvanishing criterion allows the
existing determinant estimates to go further. The arithmetic estimate
still uses class number one to control non-inert primes by norm
representations. A general argument for Landau--Siegel zeros would
require a substitute for that input.

\Needspace{16\baselineskip}
\section*{AI use}\label{sec:ai-use}
The underlying research was carried out by autonomous rStation agents
using Gemini 4 Argon and GPT-6 Astra. The author launched the agents,
which conducted the research without human intervention. The sources for
the earlier bound are
\href{https://rstation.ai/problems/tenth_discriminant/papers/2/v12}{paper 2,
version 12} and
\href{https://rstation.ai/problems/tenth_discriminant/papers/22/v1}{paper 22,
version 1}. The source for the bound $331$ is
\href{https://rstation.ai/problems/tenth_discriminant/papers/24}{rStation paper 24}.
Its calculations extend
\href{https://rstation.ai/problems/tenth_discriminant/papers/2/v13}{paper 2,
version 13}, together with the rank argument developed in
\href{https://rstation.ai/problems/tenth_discriminant/papers/22/v4}{paper 22,
version 4}. The spectral estimates originate in
\href{https://rstation.ai/problems/tenth_discriminant/papers/16/v2}{paper 16,
version 2}. 

Separate agents using GPT-6 Astra and Claude Opus 5.5 checked the original
arguments independently of the discovering agents. GPT-6 Astra, Claude
Opus 5.5 and Gemini 4 Argon were then used in preparing the present
exposition, including mathematical consultation, revision of the outline
and prose, and review of the proof. GPT-6 Astra recomputed the six numerical comparisons for this paper
and used separate programs to check the rank inequalities, the
auxiliary-prime counts and the remaining discriminants. The source
calculation at $347$ included a four-factor composite correction;
the comparisons here use only the two- and three-factor corrections.

\begin{thebibliography}{99}
\raggedright

\bibitem{baker}
Alan Baker, \emph{Linear forms in the logarithms of algebraic numbers},
Mathematika \textbf{13} (1966), no.~2, 204--216.
\href{https://doi.org/10.1112/S0025579300003971}{doi:10.1112/S0025579300003971}.

\bibitem{chung498}
Stephen Chung, \emph{A determinant bound of 498 for the tenth discriminant},
working draft, 28 October 2026.
\href{https://github.com/dualverse-ai/tenth_discriminant/blob/main/bound-498/tenth-discriminant-498.pdf}{Paper};


\bibitem{green2026}
Ben Green, \emph{On Landau--Siegel zeros},
Hexagon preprint 2610.00311, version 2, 18 October 2026.
\url{https://hexagonmath.org/pdf/2610.00311v2}.

\bibitem{heegner}
Kurt Heegner, \emph{Diophantische Analysis und Modulfunktionen},
Math. Z. \textbf{56} (1952), 227--253.
\href{https://doi.org/10.1007/BF01174749}{doi:10.1007/BF01174749}.

\bibitem{stark}
Harold M. Stark, \emph{A complete determination of the complex quadratic fields of class-number one},
Michigan Math. J. \textbf{14} (1967), 1--27.
\href{https://doi.org/10.1307/mmj/1028999653}{doi:10.1307/mmj/1028999653}.

\bibitem{taosketch}
Terence Tao, \href{https://terrytao.wordpress.com/2026/10/19/a-sketch-of-the-proof-of-non-existence-of-siegel-zeroes/}{\emph{A sketch of the proof of non-existence of Siegel zeroes}},
\emph{What's New}, 19 October 2026.

\bibitem{tao2026}
Terence Tao, \href{https://terrytao.wordpress.com/2026/10/20/an-experimental-ai-math-challenge-using-tenth-discriminant-bounds-to-explore-arguments-ruling-out-siegel-zeroes/}{\emph{An experimental AI math challenge: using tenth discriminant bounds to explore arguments ruling out Siegel zeroes}},
\emph{What's New}, 20 October 2026.

\end{thebibliography}

\end{document}
